One again obtains essentially the same results for @@M0@@ in place of @@M1@@ or @@M2@@, with the same reservation concerning the bijectivity assertion in (i). In the case where one restricts to @@M3@@, this bijectivity assertion is valid when @@M4@@ is noetherian and \emph{normal}, for in this case, for every open subset @@M5@@ of an affine space over @@M6@@, one sees immediately that @@M7@@ is surjective (EGA @@M8@@ 53, 21.4.13).

\sgaStableTarget{sga6.e10.bibliography}\subsection*{BIBLIOGRAPHY}
\begin{enumerate}[label={[\arabic*]}]
\item Dickson, E. Linear groups, Dover, New York, 1958.
\item Grothendieck, A. TDTE, V, VI; Séminaire Bourbaki 232, 236.
\item Kleiman, S. Towards a numerical theory of ampleness, Ann. of Math. 84 (1966), pp. 293-344.
\item Murre, J. On contravariant functors from the category of préschemes over a field into the category of abelian groups, Publications Mathématiques de l'IHES, 23, 1964.
\item Chevalley, C. Les classes d'équivalence rationnelles, I et II, Séminaire à l'ENS 1958, exposés 2 et 3.
\item S. Abhyankar, Resolution of singularities of arithmetical surfaces, in: Arithmetical Algebraic Geometry, Purdue University, Harper--Row, New York 1965.
\end{enumerate}

\providecommand{\Isom}{\operatorname{Isom}}
\providecommand{\Hom}{\operatorname{Hom}}
\providecommand{\Coker}{\operatorname{Coker}}
\providecommand{\Nil}{\operatorname{Nil}}
\providecommand{\im}{\operatorname{Im}}
\providecommand{\coker}{\operatorname{coker}}
\providecommand{\Supp}{\operatorname{Supp}}
\providecommand{\fppf}{\mathrm{fppf}}
\providecommand{\zar}{\mathrm{zar}}
\providecommand{\GL}{\operatorname{Gl}}
\providecommand{\Gr}{\operatorname{Gr}}
\providecommand{\Pic}{\operatorname{Pic}}
\providecommand{\Spec}{\operatorname{Spec}}
\providecommand{\V}{\operatorname{V}}
\providecommand{\Ass}{\operatorname{Ass}}

\providecommand{\Hom}{\operatorname{Hom}}
\providecommand{\Isom}{\operatorname{Isom}}
\providecommand{\Coker}{\operatorname{Coker}}
\providecommand{\Ass}{\operatorname{Ass}}
\providecommand{\sgacircled}[1]{\tikz[baseline=(char.base)]\node[shape=circle,draw,inner sep=.35pt](char){\mbox{\scriptsize #1}};}
\newpage
